% SEGMENT OLP-0665-S01
% Part: proof-theory
% Chapter: natural-deduction
% Section: quantifiers

\documentclass[../../../include/open-logic-section]{subfiles}

\begin{document}

\olfileid{pt}{nat}{qua}

\olsection{ஒழுங்குறு \usetoken{P}{proof} மற்றும் பதிலீடு}

\Elim{\lexists}, \Intro{\lforall} விதிகளுக்குப் பொருந்தும்
“தனிமாறி நிபந்தனைகள்” காரணமாக அளவையடை விதிகளில் சில
சிக்கல்கள் தோன்றுகின்றன. இந்த அனுமானங்களில் வரும்
!!{constant}~$c$ ஐ மீண்டும் பயன்படுத்தாமல் இருப்பதன் மூலம்
பெரும்பாலான சிக்கல்களைத் தவிர்க்கலாம். பின்வரும் வரையறையை
அறிமுகப்படுத்துவோம்.

\begin{defn}
இயல்பான வருவித்தல் !!{proof}~$\delta$ பின்வரும் இரு
நிபந்தனைகளையும் நிறைவுசெய்தால் அதை \emph{ஒழுங்குறு} என்போம்:
\begin{enumerate}
  \item எந்தத் தனிமாறி~$a$ உம் ஒன்றுக்கு மேற்பட்ட
  \Elim{\lexists} அல்லது~\Intro{\lforall} அனுமானங்களின்
  தனிமாறியாகப் பயன்படுத்தப்படாது;
  \item $a$ ஆனது \Elim{\lexists} அல்லது~\Intro{\lforall}
  அனுமானத்தின் தனிமாறியாக இருந்தால், முறையே அந்த அனுமானத்தின்
  சிறு முன்கோளுக்கு அல்லது முன்கோளுக்கு இட்டுச் செல்லும்
  உடனடி உள்-!!{proof} இல் மட்டுமே அது இடம்பெறும்.
\end{enumerate}
\end{defn}

% SEGMENT OLP-0665-S02
\begin{lem}\ollabel{lem:subst-regular}
$\delta$ என்பது $!C$ இல் முடியும் ஒழுங்குறு நிறுவலாகவும்,
$c$ என்பது $\delta$ இல் தனிமாறியாகப் பயன்படுத்தப்படாத
!!a{constant} ஆகவும், $t$ என்பது $\delta$ இன் எந்தத்
தனிமாறியையும் கொண்டிராத சொல்லாகவும் இருந்தால்,
$\delta$ இல் $c$ வரும் ஒவ்வொரு இடத்திலும் $t$ ஐப்
பதிலிடுவதால் கிடைக்கும் $\Subst{\delta}{t}{c}$ ஓர்
ஒழுங்குறு !!{proof} ஆகும். $\Subst{\delta}{t}{c}$ இன்
இறுதி-!!{formula} $\Subst{!C}{t}{c}$ ஆகும். $!A^x$
என்பது $\delta$ இன் திறந்த எடுகோள் எனில்,
$\Subst{!A}{t}{c}^x$ என்பது
$\Subst{\delta}{t}{c}$ இன் திறந்த எடுகோள் ஆகும்.
\end{lem}

\begin{proof}
  $\pheight{\delta}$ மீது தொகுத்தறிதலைப் பயன்படுத்துவோம்.
  $\pheight{\delta} = 0$ எனில், அது வெறும்
  எடுகோள்~$!A^x$ மட்டுமே. அப்போது
  $\Subst{\delta}{t}{c}$ என்பது~$\Subst{!A}{t}{c}^x$.

% SEGMENT OLP-0665-S03
  $\pheight{\delta} > 0$ எனில், $\delta$ இன் இறுதி
  அனுமான விதியைப் பொறுத்து நேர்வுகளைப் பிரிப்போம்.
  கூற்று இணைப்பிகளுக்கான விதிகளின் நேர்வுகள் எளிதானவை;
  அவற்றைப் பயிற்சிகளாக விடுகிறோம். $\delta$ அளவையடை
  அனுமானத்தில் முடியும் நேர்வுகளைப் பார்ப்போம்.
  \begin{enumerate}
    \item இறுதி அனுமானம் $\Intro{\lforall}$ எனில்,
    $\delta$ இன் வடிவம்
    \[
    \AxiomC{}
    \RightLabel{$\delta'$}
    \DeduceC{$!A(a,c)$}
    \RightLabel{\Intro{\lforall}}
    \UnaryInfC{$\lforall[x][!A(x,c)]$}
    \DisplayProof
    \]
    ஆகும். தொகுத்தறிதல் கருதுகோளால்
    $\Subst{\delta'}{t}{c}$ என்பது $!A(a,t)$ இன்
    ஒழுங்குறு !!{proof}. $a$ ஆனது~$\delta$ இன்
    தனிமாறி என்பதால் $a$ ஆனது $t$ இல் இடம்பெறாது. ஆகவே
    $\Subst{\delta}{t}{c}$ இன் இறுதி அனுமானம்
    \[
    \AxiomC{}
    \RightLabel{$\Subst{\delta'}{t}{c}$}
    \DeduceC{$!A(a,t)$}
    \RightLabel{\Intro{\lforall}}
    \UnaryInfC{$\lforall[x][!A(x,t)]$}
    \DisplayProof
    \]
    சரியானது: $t$ இல் $a$ இல்லாததால் முடிவிலும்
    எந்தத் திறந்த எடுகோளிலும் $a$ இல்லை.

% SEGMENT OLP-0665-S04
    \item இறுதி அனுமானம் $\Elim{\lforall}$ எனில்,
    $\delta$ இன் வடிவம்
    \[
    \AxiomC{}
    \RightLabel{$\delta'$}
    \DeduceC{$\lforall[x][!A(x,c)]$}
    \RightLabel{\Elim{\lforall}}
    \UnaryInfC{$!A(s(c),c)$}
    \DisplayProof
    \]
    ஆகும். தொகுத்தறிதல் கருதுகோளால்
    $\Subst{\delta'}{t}{c}$ என்பது
    $\lforall[x][!A(x,t)]$ இன் ஒழுங்குறு
    !!{proof}. (முடிவிலுள்ள சொல்~$s(c)$,
    $c$ ஐக் கொண்டிருக்கலாம்; கொண்டிருக்க வேண்டிய
    கட்டாயமில்லை.) $\Subst{\delta}{t}{c}$ இன்
    இறுதி அனுமானம்
    \[
    \AxiomC{}
    \RightLabel{$\Subst{\delta'}{t}{c}$}
    \DeduceC{$\lforall[x][!A(x,t)]$}
    \RightLabel{\Elim{\lforall}}
    \UnaryInfC{$!A(s(t),t)$}
    \DisplayProof
    \]
    சரியானது; மேலும்
    $!A(s(t),t) \ident \Subst{!A(s(c),c)}{t}{c}$.
    \Intro{\lexists} நேர்வும் இதைப் போன்றதே.

% SEGMENT OLP-0665-S05
  \item இறுதி அனுமானம் $\Elim{\lexists}$ எனில்,
    $\delta$ இன் வடிவம்
    \[
    \AxiomC{}
    \RightLabel{$\delta_1'$}
    \DeduceC{$\lexists[x][!A(x,c)]$}
    \AxiomC{$\Discharge{!A(a,c)}{x}$}
    \RightLabel{$\delta_2'$}
    \DeduceC{$!C$}
    \DischargeRule{\Elim{\lexists}}{x}
    \BinaryInfC{$!C$}
    \DisplayProof
    \]
    ஆகும். தொகுத்தறிதல் கருதுகோளால்
    $\Subst{\delta_1'}{t}{c}$ மற்றும்
    $\Subst{\delta_2'}{t}{c}$ ஆகியவை ஒழுங்குறு
    !!{proof}s. முதலாவது $\lexists[x][!A(x,t)]$
    ஐ நிறுவுகிறது; இரண்டாவது $!A(a,t)$ ஐத்
    திறந்த எடுகோளாகக் கொண்டு $\Subst{!C}{t}{c}$
    ஐ நிறுவுகிறது. $a$ ஆனது~$\delta$
    இன் தனிமாறி என்பதால் $a$ ஆனது $t$ இல் இடம்பெறாது.
    ஆகவே $\Subst{\delta}{t}{c}$ இன் இறுதி
    அனுமானம்
    \[
    \AxiomC{}
    \RightLabel{$\Subst{\delta_1'}{t}{c}$}
    \DeduceC{$\lexists[x][!A(x,t)]$}
    \AxiomC{$\Discharge{!A(a,t)}{x}$}
    \RightLabel{$\Subst{\delta_2'}{t}{c}$}
    \DeduceC{$\Subst{!C}{t}{c}$}
    \DischargeRule{\Elim{\lexists}}{x}
    \BinaryInfC{$\Subst{!C}{t}{c}$}
    \DisplayProof
    \]
    சரியானது:
    $\Subst{!A(x,t)}{a}{x} \ident
    \Subst{!A(a,c)}{t}{c}$; மூல முடிவு~$!C$
    இலும் எந்தத் திறந்த எடுகோளிலும் $a$ இல்லை.
    அந்தச் சொல்லிலும் அந்தத் தனிமாறி இல்லாததால், பதிலீட்டுக்குப்
    பின்னரும் அந்தத் தனிமாறி நிபந்தனை நிறைவுறுகிறது.
  \end{enumerate}
\end{proof}

% SEGMENT OLP-0665-S06
\begin{prop}\ollabel{prop:regular}
$\delta$ என்பது \Log{N1c} அல்லது \Log{N1i} இல்
!!a{proof} எனில், அதே கட்டமைப்பையும் (குறிப்பாக,
அதே !!{height} ஐயும்) கொண்ட, தனிமாறிகளின்
பெயர்மாற்றத்தில் மட்டுமே $\delta$ இலிருந்து
வேறுபடும் ஒழுங்குறு !!{proof}~$\delta'$ ஒன்று உள்ளது.
\end{prop}

\begin{proof}
இந்த நிறுவலின் நோக்கத்திற்காக மட்டும், $\delta$
இல் உள்ள ஒரு தனிமாறி அனுமானத்தின் தனிமாறி
வேறோர் அனுமானத்தின் தனிமாறியாகவும் இல்லாமல்,
அந்த அனுமானத்தின் உடனடி உள்-!!{proof} க்கு
வெளியே $\delta$ இல் எங்கும் இடம்பெறாமலும்
இருந்தால் அதை \emph{தூயது} என்போம்;
இல்லையெனில் \emph{தூய்மையற்றது} என்போம்.
$\delta$ இல் உள்ள தூய்மையற்ற தனிமாறி
அனுமானங்களின் எண்ணிக்கை $n$ என்க.
$n$ மீது தொகுத்தறிதலால் $\delta$ வை வேண்டிய
ஒழுங்குறு !!{proof}~$\delta'$ ஆக மாற்றலாம்
என்று காட்டுவோம். $n=0$ எனில்,
$\delta$ இன் எல்லாத் தனிமாறி அனுமானங்களும்
தூயவை; எனவே $\delta$ ஒழுங்குறு, நிறுவ
வேறொன்றுமில்லை. இல்லையெனில் மிக மேலுள்ள
தூய்மையற்ற தனிமாறி அனுமானம் ஒன்றைத் தேர்வோம்;
அது \Intro{\lforall} என்க. அதில் முடியும்
$\delta$ இன் உள்-!!{proof}~$\delta_1$ இன் வடிவம்
    \[
    \AxiomC{}
    \RightLabel{$\delta_2$}
    \DeduceC{$!A(c)$}
    \RightLabel{\Intro{\lforall}}
    \UnaryInfC{$\lforall[x][!A(x)]$}
    \DisplayProof
    \]
ஆகும்.

% SEGMENT OLP-0665-S07
$c$ தனிமாறி என்பதால்
$\lforall[x][!A(x)]$ இலோ திறந்த எடுகோளிலோ
இடம்பெறாது. $\delta_2$ இல் தூய்மையற்ற
தனிமாறி அனுமானம் இல்லை; எனவே அது ஒழுங்குறு.
$\delta$ இல் இடம்பெறாத !!a{constant}~$a$
ஒன்றைத் தேர்வோம். \olref{lem:subst-regular} படி,
$\Subst{\delta_2}{a}{c}$ என்பது $!A(a)$ இன்
ஒழுங்குறு நிறுவல். தனிமாறி நிபந்தனையால்
$\delta_2$ இன் எந்தத் திறந்த எடுகோளிலும்
$c$ இல்லை; ஆகவே
$\Subst{\delta_2}{a}{c}$ இன் எந்தத் திறந்த
எடுகோளிலும் $a$ இல்லை. எனவே
$\Intro{\lforall}$ ஐப் பயன்படுத்தலாம்.
$\delta$ இல் $\delta_1$ க்குப் பதிலாகப் பின்வரும்
நிறுவல்~$\delta_1'$ ஐ வைத்தால்,
    \[
    \AxiomC{}
    \RightLabel{$\Subst{\delta_2}{a}{c}$}
    \DeduceC{$!A(a)$}
    \RightLabel{\Intro{\lforall}}
    \UnaryInfC{$\lforall[x][!A(x)]$}
    \DisplayProof
    \]
ஒரு தூய்மையற்ற தனிமாறி அனுமானத்தைத்
தூய அனுமானத்தால் மாற்றியிருக்கிறோம்;
வேறுபாடு தனிமாறி $c$ க்குப் பதிலாக~$a$
வைத்திருப்பதே. விளைவாகக் கிடைக்கும் நிறுவலில்
தூய்மையற்ற தனிமாறி அனுமானங்கள் அதிகபட்சம்
$n-1$ மட்டுமே; தொகுத்தறிதல் கருதுகோள்
முடிவைத் தருகிறது.
\end{proof}

% SEGMENT OLP-0665-S08
\begin{prob}
\olref{prop:regular} இன் நிறுவலில்,
மிக மேலுள்ள தூய்மையற்ற தனிமாறி அனுமானம்
\Elim{\lexists} ஆக இருக்கும் நேர்வை விவரிக்கவும்.
\end{prob}

இப்போது நிறுவிய முன்மொழிவை இனி வெளிப்படையாக
மேற்கோள் காட்ட மாட்டோம். ஆனால் நாம் கருதும்
!!{proof}s அனைத்தும் ஒழுங்குறு என்று கொள்வோம்:
அப்படி இல்லாதவற்றின் தனிமாறிகளைப்
பெயர்மாற்றி ஒழுங்குறச் செய்யலாம்.

\Olref{lem:subst-regular} உம்
\olref{prop:regular} உம் \Log{N2c} மற்றும்~
\Log{N2i} க்கும் பொருந்தும். அவற்றின்
நிறுவல்களைப் பயிற்சிகளாக விடுகிறோம்.

\end{document}
